% Section fragment. Requires amsmath, amssymb, amsthm and hyperref.
\section{PA194GH: explicit generalization and cached-head halting}
\label{sec:gh}

PA194GH combines an AX-supplied generalization binder with no-push modus
ponens and a Shift instruction that tests only its cached formula.
This combination is a different command interface from the earlier PA194G
and PA194H interfaces. Its state count is still 194. The five literal
changes from PA194G are used by the recurring Shift controller; this is not
renaming, reflection, a different initial state, or a one-time padding change.
The proof below establishes a common numerical-calculus interface. It does
not claim that the repository's unfinished calibration of that calculus
and physical machine to PA has become a completed Lean/iPRA theorem.

\subsection{Exact command semantics}

Use Cantor pairing $p(a,b)=(a+b)(a+b+1)/2+b$, total stack constructors
$[a;w]=p(a,w)$, and total projections $\mathrm{hd},\mathrm{tl}$.
Thus $[0;0]=0$ and both projections of zero are zero. Put
$I(a,b)=4p(a,b)$ and $U(a,b)=4p(b,a)+2$. A store is $(t,w,d,h)$.
The command semantics is
\[
 S(t,w,d,h)=\begin{cases}\mathsf{halt}&h=1,\\
 (h,\mathrm{tl}(w),d,\mathrm{hd}(w))&h\ne1,
 \end{cases}
\]
and, for $b=d+x$,
\[
 \mathsf{MP}_{x}(t,w,d,h)=\begin{cases}(b,w,b,0)&h=I(t,b),\\
 (0,w,b,0)&h\ne I(t,b).
 \end{cases}
\]
In particular, MP pushes no extra zero, and successful MP leaves its
conclusion in both the visible top and the offset.

For AX let $D=d+x$, $H=h+y$, $A=c$. It leaves the inherited guarded gate
$K(H,D,A)$ as top, clears both caches, and pushes, in stack-top order,
\[
 q_1=\Gamma(A),\quad (q_2,\ldots,q_7)=(L_6,\ldots,L_1),\quad
 q_8=\rho_8(D,A,H,d),\quad q_9=U(A,h),\quad t.
\]
Here $\Gamma(A)=I(A,U(A,A))$. The $L_i$, the guard and gate, and the
dominated implicational head are exactly those emitted by PA194G; its
exchanged distribution orientation is retained. The last universal head
contains the \emph{incoming} $h$, not $H=h+y$. This distinction is required
for soundness: $h$ is a theorem under the store invariant; $h+y$ need not be.
The binder is the explicit operand $c$, independent of the incoming offset.
Every $q_j$ and the gate differs from one: $q_1,\ldots,q_8,K$ have residue
zero modulo four, and $q_9$ has residue two. A failed guard produces zero.

\subsection{All-store identities and zero aliases}

Words execute left to right; exponentiation of a word denotes repetition.
Define
\[
 \mathsf{Push}=\mathsf{AX}_{0,0,0};S^9;\mathsf{MP}_0,
 \quad \mathsf{Pop}=\mathsf{AX}_{0,0,0};S^{11},
 \quad \mathsf{Discard}=\mathsf{MP}_0.
\]
The following are equalities of total numerical encodings, not arguments
about a positive stack length:
\begin{align*}
 \mathsf{AX}_{c,x,y};S^9(t,w,d,h)&=(q_8,[t;w],0,U(c,h)),\\
 \mathsf{AX}_{c,x,y};S^{10}(t,w,d,h)&=(U(c,h),w,0,t),\\
 \mathsf{Push}(t,w,d,h)&=(0,[t;w],0,0),\\
 \mathsf{Discard}(t,w,0,h)&=(0,w,0,0),\\
 \mathsf{Pop}(t,[a;w],d,h)&=
 \begin{cases}\mathsf{halt}&t=1,\\(t,w,0,a)&t\ne1.\end{cases}
\end{align*}
The first ten Shifts test only the cleared cache and the generated heads,
not the saved arbitrary top. In Push, the last MP cannot match a
residue-zero implication against the residue-two universal head. In
Discard its proposed conclusion is zero, so success and failure have the
same result. The eleventh Shift in Pop tests the saved old top and is its
only possible halt. Empty-stack cases follow directly from $[0;0]=0$;
for example Push of the all-zero store returns that same numeric store.
Pop must never be described as total returning execution when $t=1$.

Generalization has the shorter compiler
\[
 \mathsf{Gen}_i=\mathsf{Push};S;\mathsf{AX}_{i,0,0};S^{10}.
\]
For every $a,w,i$, including $a=1$ and $i=0$, it returns
$(U(i,a),w,0,0)$ from $(a,w,0,0)$. After Push the single Shift tests zero
and puts $a$ in the cache. AX consumes that cache before another Shift
can inspect it. PA194H requires an additional offset-setting MP here;
PA194GH does not.

For $1\le j\le7$, the selector
\[
 \mathsf{Discard};\mathsf{AX}_{A,D,H};S^{j+1};\mathsf{Pop}^{9-j}
\]
returns $(q_j,w,0,0)$ from every $(t,w,0,0)$. The gate selector is
$\mathsf{Discard};\mathsf{AX}_{c,b,a};\mathsf{Pop}^{10}$.
For $u\le a$, write $\mathsf{Offset}_u=\mathsf{MP}_u$; the selector
\[
 \mathsf{Discard};\mathsf{Offset}_u;
 \mathsf{AX}_{b,a-u,c};S^9;\mathsf{Pop}
\]
returns $\rho_8(a,b,c,u)$ with the original underlying stack and cleared
caches. All selected formulas in these words have residue zero, so their
Pops are safe. The saved zero is consumed last.

\subsection{The MP compiler and its reserved zero}

Put $\mathsf{MPSuffix}_b=S;S;\mathsf{MP}_b;\mathsf{Pop}$. Then
\[
 \mathsf{MPSuffix}_b(0,[a,I(a,b),0;w],0,0)=
 \begin{cases}\mathsf{halt}&a=1\text{ or }b=1,\\
 (b,w,0,0)&a\ne1\text{ and }b\ne1.
 \end{cases}
\]
The second Shift detects $a=1$. Otherwise MP sees minor $a$ and cached
major $I(a,b)$, succeeds, and leaves $(b,[0;w],b,0)$. The final Pop
clears the offset and consumes the reserved zero, or detects $b=1$.
Without this zero it would remove part of the arbitrary underlying stack.
The case $a=b=w=0$ uses the literal identities $I(0,0)=0$ and $[0;0]=0$,
and is included in the calculation.

For a proof tree $\pi$ in the common numerical calculus define a word
$W_\pi$ recursively. Use the displayed selectors for axiom instances and
$W_{\pi_0};\mathsf{Gen}_i$ at a generalization node. At an MP node,
where $\pi_0$ concludes $a$ and $\pi_1$ concludes $I(a,b)$, use
\[
 \mathsf{Discard};\mathsf{Push};
 W_{\pi_1};\mathsf{Push};W_{\pi_0};\mathsf{Push};\mathsf{MPSuffix}_b.
\]
For all $t,w$, this either halts early or returns
$(\operatorname{concl}(\pi),w,0,0)$ from $(t,w,0,0)$.
This follows by proof-tree induction using the exact word identities.
Each recursive word leaves its arbitrary old stack unchanged, so the
reserved zero survives both premise computations.

With $n$ proof-tree nodes, $|W_\pi|\le122n$: the longest axiom selector
is the 122-command gate selector; generalization adds 23 commands; an
MP node adds 49 commands in addition to its two subwords. Consequently,
if $\pi$ proves one, the word
\[
 V_\pi=\mathsf{AX}_{0,0,0};W_\pi;\mathsf{Pop}
\]
halts from every incoming store and has length at most $122n+13$.
Its first AX clears dirty caches, and the final Pop detects the proven
one before consulting arbitrary old stack data. A final lone Shift would
be wrong for this interface: $S(1,0,0,0)=(0,0,0,0)$.

\subsection{Fair scheduling and a primitive-recursive witness}

The carried-store scheduler is unchanged. For its command codes write
\[
 B_S(z)=p(0,z),\quad B_{\mathsf{MP}_x}(z)=p(1,p(x,z)),\quad
 B_{\mathsf{AX}_{c,x,y}}(z)=p(c+2,p(y,p(x,z))).
\]
The continuation encoding is $E(\epsilon)=0$ and
$E(c::w)=1+B_c(E(w))$. If $z=E(V_\pi)$ then the fresh-candidate codes
\[
 n_k=p(k+2,p(0,p(0,z)))
\]
are strictly increasing, cofinal, and execute
$\mathsf{AX}_{k,0,0};V_\pi$. The prefix may leave any store because
$V_\pi$ halts from every store. Given a current fresh counter $N$, choose
$k=N$; then $n_k>N$. The current finite continuation is first completed,
unless it halts. For a positive continuation $l$ the next residual is at
most $l-1$, and a fresh candidate $j$ therefore has at most $j+1$ commands.
For example
\[
 L+(n_k+2)^2+2(n_k+2)
\]
is a safe bound on remaining abstract scheduler command/boundary events
through this chosen candidate, where $L$ is the current continuation.
The bound is deliberately generous and primitive recursive. It proves
existence of a finite halting witness without requiring reset stores or
standardness of the natural-number model.

Conversely annotate $t$, $h$, and the live canonical Cantor-stack entries
with numerical-calculus derivations. The all-zero initial store has the
fixed proof of zero (the all-zero dominated-head instance); it also
supplies each virtual zero tail. AX appends proofs of its axiom heads and
uses theorem-generalization on the incoming proof of $h$ to obtain
$U(c,h)$. MP uses the annotated premises on success, and the fixed zero
proof on failure. Shift moves annotations and halts only with an
annotated derivation of one. Pairing aliases are handled by canonical
expansion and the explicit zero proof, not by assuming every push creates
a nonempty stack cell. Thus a finite scheduler halting trace produces a
finite proof of one.

Both directions are witness-producing. Encode proof trees with proper
subtrees smaller than their root code, and use primitive-recursive
course-of-values recursion for $W_\pi$. The soundness annotation is a
primitive-recursive fold over a finite trace. Proof validity, store
annotation validity, decoded list bounds, and bounded scheduler runs all
have primitive-recursive characteristic functions. The inductions just
used can therefore be expressed as quantifier-free induction with those
functions in a PRA language. All branches test decidable numerical
relations. For the all-store compilation lemma, do not induct on an
unbounded universally quantified store assertion. Fix the proof code and
initial $t,w$ as free parameters; compute the finite word, its full
length-bounded execution with an absorbing halt flag, and the finite tree
of recursive subword intervals. Every actual recursive call has a store
appearing in that generated trace. Its entries are bounded by the
primitive-recursively computed maximum of the trace encoding. At each
subtree rank, use the bounded conjunction of the desired identities for
these finitely many actual call records. Its characteristic function is
primitive recursive, so rank induction is genuinely quantifier-free.
The preceding all-store statement follows because the initially fixed
$t,w$ were arbitrary; no unbounded universal statement was used as the
induction formula. In particular, this new command-level proof has a constructive
primitive-recursive formulation; a bare invocation of induction available
only in PA is not being substituted for it.

The source-to-tape step requires the explicit loop realizations of the
arithmetic commands and their invariant. Its local tape components have
positive affine durations in the arbitrary register values and header
gap. Summing these durations along a primitive-recursively generated
source trace gives a primitive-recursive physical time bound. The supplied
literal-table checker establishes the relevant local identities and list
induction certificates; formal internalization of their checker and the
full inherited command implementation remains a separate proof-assistant
obligation. Accordingly the unconditional new result at this level is the
witness equivalence between the displayed abstract scheduler and the
common numerical calculus, with an explicit checked physical realization
as documented next.

\subsection{Physical realization and exact scope}

The AX grammar is PA194G's frozen 190-token grammar. Its 67 reader controls
are combined with 134 fixed controls, then seven disjoint contextual pairs
are folded: $134+67-7=194$. The no-push source subtree remains
\texttt{node163=(node6,node54)}. The Shift subtree becomes
\[
 \texttt{node170=((node36,(node167,node145)),node57)},
\]
which is PA194H's clear-old-top, test-cached-head implementation.
The five fixed row edits are
\[
\begin{array}{c|c|c}
(q,b)&\hbox{PA194G action}&\hbox{PA194GH action}\\\hline
(38,0)&(0,R,28)&(0,R,11)\\
(38,1)&(1,R,20)&(1,R,122)\\
(214,0)&(0,R,188)&(0,R,312)\\
(214,1)&(1,R,312)&(1,R,20)\\
(312,0)&(0,R,11)&(0,R,188).
\end{array}
\]
The entire resolved source has 3,653 primitive sites and 4,754 outcomes,
including both decrement branches. All 28 two-cell fold identities hold
simultaneously. The constructor checks every outcome with arbitrary
register symbols, and checks equality of the complete blank-start tape
with the intended frame after exactly 57,922 transitions. The separate
literal checker rederives the four unary-list sweep lemmas from the
actual candidate rows and does not assume the builder's arithmetic-box
oracle. These are universal symbolic certificates, not a sample run over
small register values, and not a completed Lean proof.

The correct endpoint to report is relative. Once the shared physical
implementation/calculus bridges are internalized in the specified weak
arithmetic, the witness transforms give halting equivalence of PA194GH
and the primary calibrated calculus (and hence of their consistency
negations in iPRA). If only the older PA calibration is assumed, the
resulting absolute consistency theorem has that older PA scope.
An unconditional machine-specific iPRA equivalence with
$\operatorname{Con}(\mathrm{PA})$ must still discharge the repository's
open physical and logical-calibration obligations. Changing five rows
and completing the new command proof cannot by itself discharge them.

\subsection{PA194GHS: a second physical implementation}

PA194GHS implements the same GH commands with the symmetric comparison
method inherited from PA194HS. It is not a second new logical calculus.
The nine source registers are $(N,T,W,A,D,H,X,Y,L)$. The GH Shift path
clears the old $T$, copies $H$ into $T$ and $X$, restores $X$ to $H$,
increments $X$, and enters the destructive comparison with
$(T,H,X)=(h,h,1)$. GHS omits that restore loop and increments $H$ instead,
entering with $(T,H,X)=(h,1,h)$. Equality is symmetric, and the unequal
branch clears both compared registers before the same Cantor-stack
continuation. Hence both programs have exactly the GH Shift semantics.
The no-push MP copy tail sends $(T,D,H,X)=(0,b,0,0)$ to $(b,b,0,0)$
in $7b+2$ primitive steps; it is unchanged in this companion.

The literal source changes are
\[
\begin{split}
\texttt{node170}&=\texttt{((node36,(node167,node145)),node56)},\\
\texttt{node171}&=\texttt{(node170,(inc5,continue3))}.
\end{split}
\]
Before folding, the four row settings are
$(214,1)\mapsto(1,R,30)$, $(310,1)\mapsto(1,R,330)$,
$(330,0)\mapsto(0,R,63)$, and $(330,1)\mapsto(1,L,73)$.
The additional contextual pair is $41\mapsto330$, with folded actions
$0L31$ and $1L46$. This gives $135+67-8=194$ working states.
The full final simultaneous table passes all 32 local fold cases.
There are 3,650 source sites and 4,750 outcomes; the constructor checks
all outcomes with arbitrary register symbols, and startup reaches the
same whole-tape frame after 57,768 transitions. The full command proof
above applies unchanged. Its independent literal-table verification is
recorded separately from GH's, so the smaller source count is not itself
used as a correctness argument.
